\Askhsh{24131}{4}
\begin{ylh}{1}
    \item 1.2 Συναρτήσεις
    \item 1.3 Μονότονες συναρτήσεις - Αντίστροφη συνάρτηση
    \item 3.5 Η συνάρτηση [ορισμένο ολοκλήρωμα της f από α έως χ]
    \item 3.7 Εμβαδόν επιπέδου χωρίου.
\end{ylh}
\begin{wrapfigure}[15]{r}{0.6\textwidth}
\begin{tikzpicture}
\begin{axis}[
    axis lines = middle,
    xlabel = {$x$},
    ylabel = {$y$},
    xmin = -0.7, xmax = 7.5,
    ymin = -0.7, ymax = 4.5,
    xtick = {-0.5,0,0.5,1,1.5,2,2.5,3,3.5,4,4.5,5,5.5,6,6.5,7},
    ytick = {-0.5,0,0.5,1,1.5,2,2.5,3,3.5,4},
    grid = both,
    grid style = {gray!30},
    samples = 200,
    width = \textwidth * 0.6, height = \textwidth * 0.45,
    tick label style = {font=\scriptsize},
    every axis x label/.style = {at={(current axis.right of origin)}, anchor=north west},
    every axis y label/.style = {at={(current axis.above origin)}, anchor=south},
]
    \addplot[thick, gray, domain=-0.7:4.5] {x};
    \node[black, below right] at (axis cs:3,3) {$y=x$};

    \addplot[very thick, maincolor, domain=0:7] { (sqrt(x)-1)/(sqrt(x)+2) };
    \node[maincolor, above right] at (axis cs:5.5,0.25) {$C_1$};

    \addplot[very thick, maincolor, domain=-0.49:0.35] { ((2*x+1)/(1-x))^2 };
    \node[maincolor, above left] at (axis cs:0.25,3.7) {$C_2$};
\end{axis}
\end{tikzpicture}
\end{wrapfigure}
Δίνεται η γνησίως αύξουσα συνάρτηση $f$, με τύπο $f(x) = \frac{\sqrt{x}-1}{\sqrt{x}+2}$, $x \ge 0$.

\begin{alist}
\item Να βρείτε το σύνολο τιμών της $f$. \monades{07}
\item Να βρείτε την αντίστροφη της $f$. \monades{07}
\item Στο παρακάτω σχήμα δίνονται οι καμπύλες $C_1, C_2$. Με δεδομένα ότι
    \begin{itemize}
        \item η μία από τις δύο καμπύλες αντιστοιχεί στην γραφική παράσταση της $f$ και η άλλη στην γραφική παράσταση της $f^{-1}$,
        \item $\int_{-\frac{1}{2}}^{0} f^{-1}(x) dx = \alpha$
    \end{itemize}
    Να βρείτε:
    \begin{rlist}
    \item Ποια καμπύλη παριστάνει την γραφική παράσταση της $f$ και ποια την γραφική παράσταση της $f^{-1}$, \monades{04}
    \item Το πρόσημο του $\alpha$ καθώς και το ολοκλήρωμα $I = \int_{0}^{1} f(x)dx$ συναρτήσει του $\alpha$. \monades{07}
    \end{rlist}
\end{alist}